\documentclass{article}
\usepackage[a4paper,margin=1in]{geometry}
\usepackage{fancyhdr}

\usepackage[yyyymmdd]{datetime}
\renewcommand{\dateseparator}{-}

\usepackage{amsmath}
\usepackage{amssymb}
\usepackage{amsthm}

\usepackage{titlesec}

\renewcommand*\contentsname{Inhalt}
\renewcommand*\proofname{Beweis}

\titleformat{\section}
{\normalfont\large\bfseries}{\thesection.}{1em}{}

\lhead{Emma Bach}
\rhead{\today}

\pagestyle{fancy}

%AMS packages - Symbols, "Theorem" and "Proof" Environments
\usepackage{amsmath}
\usepackage{amssymb}
\usepackage{amsthm}

\usepackage{epigraph}
\usepackage{physics}

\usepackage{multirow}

\usepackage[yyyymmdd]{datetime}
\renewcommand{\dateseparator}{-}


% Tikz and Graphics
\usepackage{tikz}
\usepackage{tikz-cd}
\usepackage{xcolor}
\usepackage{quiver}

% Nicer Underlining
\usepackage{contour}
\usepackage[normalem]{ulem}

% Multiple Text Columns
\usepackage{multicol}

% Margin Notes
\usepackage{sidenotes}

% FloatBarrier
\usepackage{placeins} 


%--- RENEWED COMMANDS ---

% Variant Greek Letters
\renewcommand\epsilon{\varepsilon}
\renewcommand{\varphi}

% Angle Symbol
\renewcommand{\angle}[2]{\measuredangle(#1,#2)}

\DeclareMathAlphabet{\mathbbold}{U}{bbold}{m}{n}
\newcommand*{\boldone}{\mathbbold{1}}

\usepackage[tight]{minitoc}
\usepackage{etoolbox}

%
\fancypagestyle{noheader}{
	\fancyhead{}
	\renewcommand{\headrulewidth}{0pt}
}
\pretocmd{\parttoc}{\begingroup\pagestyle{noheader}}{}{}
\apptocmd{\parttoc}{\endgroup}{}{}

% Titles and stuff set in the document itself

%--- TEXT FORMATING ---

\newcommand{\+}{\quad+}

\newcommand{\separator}{\vspace{-8pt}\hrule\vspace{8pt}}

% nice underlining
\renewcommand{\ULdepth}{1.6pt}
\contourlength{0.8pt}
\newcommand{\ul}[1]{%
	\uline{\phantom{#1}}%
	\llap{\contour{white}{#1}}%
}

% no indent
\setlength{\parindent}{0pt}

% shorthand
\newcommand{\tbf}[1]{\textbf{#1}}
\newcommand{\tn}[1]{\textnormal{#1}}
\newcommand{\tit}[1]{\textit{#1}}
\newcommand{\ttt}[1]{\texttt{#1}}
\newcommand{\tib}[1]{\tbf{\emph{#1}}}

\newcommand{\on}{\operatorname}

\newcommand{\eps}{\epsilon}
\newcommand{\lb}{\lambda}
\newcommand{\tl}{\tilde}

% make \textbf automatically include math
\let\oldtextbf=\textbf
\renewcommand\textbf[1]{{\boldmath\oldtextbf{#1}}}

\newcommand{\mc}[1]{\mathcal{#1}}
\newcommand{\ol}[1]{\overline{{#1}}}

% Underlined, bold, non-cursive Theorem Name
\newcommand{\theoremname}[1]{\textnormal{\tbf{(#1):}}}

% Starts a new paragraph without indentation
% and with an empty line between paragraphs
\newcommand*{\newpar}{\par\vspace{\baselineskip}\noindent}

% Make inf height match sup height
\renewcommand{\inf}{\mathop{\mathrm{inf}\vphantom{\mathrm{sup}}}}

% Make tildes more readable
\renewcommand{\tilde}{\widetilde}

%--- RELATION SYMBOLS AND OPERATORS ---
\newcommand{\surj}{\twoheadrightarrow}
\newcommand{\inj}{\hookrightarrow}
\newcommand{\iso}{\overset{\sim}{\rightarrow}}
\newcommand{\symdiff}{\vartriangle}
\newcommand{\trans}{\twoheadrightarrow}
\newcommand{\tensor}{\otimes}
\newcommand{\Tensor}{\bigotimes}
\newcommand{\bigmid}{\Big |}
\renewcommand{\mid}{~\middle|~}

\newcommand{\subseteqop}{~\mathring \subseteq~}
\newcommand{\subseteqcl}{~\bar \subseteq~ \vphantom{\mathring \subseteq}}

\newcommand{\supseteqop}{~\mathring \supseteq~}
\newcommand{\supseteqcl}{~\bar \supseteq~ \vphantom{\mathring \subseteq}}

\newcommand{\opensub}[2]{\underset{\subseteqop #2}{#1}}
\newcommand{\opensubil}[2]{{#1}_{\subseteqop #2}}

%--- STUFF IN FANCY BRACKETS ---
\newcommand{\stack}[2]{\begin{array}{c} #1\\#2\end{array}}
\newcommand{\trinom}[3]{\begin{pmatrix} #1\\#2\\#3\end{pmatrix}}
\newcommand{\scalar}[2]{\left\langle #1, #2 \right\rangle}
\newcommand{\angles}[1]{\left\langle #1 \right\rangle}
\newcommand{\lr}{\qty}


%--- LETTERS ---
\newcommand{\dmu}{\ d\mu}
\newcommand{\dist}{\textnormal{dist}}
\newcommand{\spt}{\textnormal{spt}}
\newcommand{\del}{\partial}

\newcommand{\pd}[2]{\frac{\del #1}{\del #2}}

% mathbb
\newcommand{\bk}{\Bbbk}
\newcommand{\bB}{\mathbb{B}}
\newcommand{\bC}{\mathbb{C}}
\newcommand{\bE}{\mathbb{E}}
\newcommand{\bF}{\mathbb{F}}
\newcommand{\bH}{\mathbb{H}}
\newcommand{\bN}{\mathbb{N}}
\newcommand{\bP}{\mathbb{P}}
\newcommand{\bQ}{\mathbb{Q}}
\newcommand{\bR}{\mathbb{R}}
\newcommand{\bS}{\mathbb{S}}
\newcommand{\bV}{\mathbb{V}}
\newcommand{\bZ}{\mathbb{Z}}

% mathcal
\newcommand{\cA}{\mathcal{A}}
\newcommand{\cB}{\mathcal{B}}
\newcommand{\cC}{\mathcal{C}}
\newcommand{\cD}{\mathcal{D}}
\newcommand{\cE}{\mathcal{E}}
\newcommand{\cF}{\mathcal{F}}
\newcommand{\cH}{\mathcal{H}}
\newcommand{\cI}{\mathcal{I}}
\newcommand{\cL}{\mathcal{L}}
\newcommand{\cM}{\mathcal{M}}
\newcommand{\cN}{\mathcal{N}}
\newcommand{\cO}{\mathcal{O}}
\newcommand{\cP}{\mathcal{P}}
\newcommand{\cQ}{\mathcal{Q}}
\newcommand{\cR}{\mathcal{R}}
\newcommand{\cS}{\mathcal{S}}
\newcommand{\cT}{\mathcal{T}}
\newcommand{\cU}{\mathcal{U}}
\newcommand{\cV}{\mathcal{V}}
\newcommand{\cW}{\mathcal{W}}
\newcommand{\cX}{\mathcal{X}}
\newcommand{\cZ}{\mathcal{Z}}

%mathfrac
\newcommand{\fa}{\mathfrak{a}}


% real numbers in fancy costumes
\newcommand{\barR}{\ol{\mathbb{R}}}
\newcommand{\bRnn}{\mathbb{R}^{n \times n}}
\newcommand{\bRmn}{\mathbb{R}^{m \times n}}
\newcommand{\bRnm}{\mathbb{R}^{n \times m}}

% ol
\newcommand{\oz}{{\ol{z}}}
\newcommand{\ow}{{\ol{w}}}

% vectors
\newcommand{\vzero}{\kern-1.2pt\vec{\kern1.2pt 0}} 
\renewcommand{\va}{\vec{a}}
\newcommand{\vA}{\vec{A}}
\renewcommand{\vb}{\vec{b}}
\newcommand{\vB}{\vec{B}}
\newcommand{\vc}{\vec{c}}
\newcommand{\ve}{\vec{e}}
\newcommand{\vE}{\vec{E}}
\newcommand{\vF}{\vec{F}}
\newcommand{\vg}{\vec{g}}
\newcommand{\vh}{\vec{h}}
\newcommand{\vH}{\vec{H}}
\newcommand{\vp}{\vec{p}}
\newcommand{\vr}{\vec{r}}
\newcommand{\vs}{\vec{s}}
\newcommand{\vt}{\vec{t}}
\newcommand{\vT}{\vec{T}}
\renewcommand{\vu}{\vec{u}}
\newcommand{\vU}{\vec{U}}
\newcommand{\vv}{\vec{v}}
\newcommand{\vV}{\vec{V}}
\newcommand{\vw}{\vec{w}}
\newcommand{\vW}{\vec{W}}
\newcommand{\vx}{\vec{x}}
\newcommand{\vy}{\vec{y}}
\newcommand{\vz}{\vec{z}}

\newcommand{\veta}{\vec{\eta}}
\renewcommand{\grad}{\vec{\nabla}}

% letters with tilde
\newcommand{\tx}{\tilde x}
\newcommand{\ty}{\tilde y}
\newcommand{\tb}{\tilde b}
\newcommand{\tPhi}{\tilde \Phi}

% bold letters
\newcommand{\tbA}{\mathbf{A}}
\newcommand{\tbB}{\mathbf{B}}
\newcommand{\tbC}{\mathbf{C}}
\newcommand{\tbD}{\mathbf{D}}
\newcommand{\tbE}{\mathbf{E}}
\newcommand{\tbY}{\mathbf{Y}}
\newcommand{\tbZ}{\mathbf{Z}}

% sequences
\newcommand{\an}{(a_n)_{n \in \bN}}
\newcommand{\bn}{(b_n)_{n \in \bN}}
\newcommand{\sn}{(s_n)_{n \in \bN}}
\newcommand{\iinI}{_{i \in I}}
\newcommand{\iinN}{_{i \in \bN}}

% special functions
\newcommand{\at}{\on{at}}
\newcommand{\cod}{\on{cod}}
\newcommand{\codim}{\on{codim}}
\newcommand{\dom}{\on{dom}}
\newcommand{\ddt}{\frac{d}{dt}}
\newcommand{\ddtz}{\frac{d}{dt}|_{t = 0}}
\newcommand{\ddx}{\frac{d}{dx}}
\renewcommand{\ev}{\on{ev}}
\newcommand{\id}{\on{id}}
\newcommand{\Id}{\on{Id}}
\newcommand{\im}{\on{im}}
\newcommand{\inv}{\on{inv}}
\newcommand{\mor}{\on{mor}}
\newcommand{\ob}{\on{ob}}
\newcommand{\ord}{\on{ord}}
\newcommand{\rang}{\on{rang}}
\renewcommand{\tr}{\on{tr}}
\newcommand{\vol}{\on{vol}}
\newcommand{\supp}{\on{supp}}
\newcommand{\cond}{\on{cond}}
\newcommand{\sgn}{\on{sgn}}
\renewcommand{\char}{\on{char}}
\newcommand{\Frac}{\on{Frac}}
\newcommand{\card}{\on{card}}
\newcommand{\intr}{\on{int}}
\newcommand{\extr}{\on{ext}}
\newcommand{\diam}{\on{diam}}

% induced set systems
\newcommand{\sigalg}[1]{\angles{#1}_\sigma}
\newcommand{\boolalg}[1]{\angles{#1}_\cB}
\newcommand{\mono}[1]{\angles{#1}_\cM}
\newcommand{\monotone}[1]{\angles{#1}_\cM}

% groups, sets and categories
\newcommand{\Moeb}{\tn{Möb}}
\newcommand{\GL}{\tn{GL}}
\newcommand{\PGL}{\tn{PGL}}
\newcommand{\SO}{\tn{SO}}
\newcommand{\SL}{\tn{SL}}
\newcommand{\PSL}{\tn{PSL}}
\newcommand{\Hess}[1]{\tn{Hess}(#1)}
\renewcommand{\div}{\tn{div}}

\newcommand{\rel}{\operatorname{rel}}

\newcommand{\Field}{\ttt{Field}}
\newcommand{\Group}{\ttt{Group}}
\newcommand{\Mat}{\ttt{Mat}}
\newcommand{\Mod}{\ttt{Mod}}
\newcommand{\Poset}{\ttt{Poset}}
\newcommand{\Ring}{\ttt{Ring}}
\newcommand{\Set}{\ttt{Set}}
\newcommand{\Top}{\ttt{Top}}
\renewcommand{\Vec}{\ttt{Vec}}

\newcommand{\Aut}{\tn{Aut}}
\newcommand{\Eig}{\tn{Eig}}
\newcommand{\Ens}{\tn{Ens}}
\newcommand{\End}{\tn{End}}
\newcommand{\Hom}{\tn{Hom}}
\newcommand{\Mag}{\tn{Mag}}
\newcommand{\Mon}{\tn{Mon}}


\newcommand{\Matnn}{\ttt{Mat}(n \times n)}
\newcommand{\MatBB}[3]{\ttt{Mat}^{#1}_{#2}\left(#3\right)}

% --- THEOREM AND PROOF TYPES ---
% \newtheorem{codename}{printedname}[countedwith]
\theoremstyle{definition}
\newtheorem{lemma}{Lemma}[section]
\theoremstyle{theorem}
\newtheorem{exercise}[lemma]{Exercise}

\theoremstyle{definition}
\newtheorem{theorem}[lemma]{Theorem}
\newtheorem{proposition}[lemma]{Proposition}
\newtheorem{corollary}[lemma]{Corollary}
\newtheorem{definition}[lemma]{Definition}
\newtheorem{example}[lemma]{Example}
\newtheorem{counterexample}[lemma]{Counterexample}
\newtheorem{anmerkung}[lemma]{Comment}

% --- CUSTOM ENVIRONMENTS ---

% sketchy proofs marked with untrustworthy QED symbol
\newenvironment{proofsketch}{\begin{proof}[Beweisskizze]\renewcommand*{\qedsymbol}{\("\square"\)}}{\end{proof}}

% align with exactly one number for labeling
\NewDocumentEnvironment{nalign}{}{\equation\aligned}{\endaligned\endequation}

\begin{document}


\title{\Huge\vspace{-1.5cm}\textbf{Lineare Algebra I im WS 26/27\\Musterlösung 04}}
\author{Emma Bach}

\maketitle

\section*{Übung 4.1}
\separator

\tit{Man zeige: Gegeben ein surjektiver Vektorraumhomomorphismus \emph{(also eine surjektive lineare Abbildung)} $g :U \surj V$ und eine Abbildung $f : V \to W$ in einen weiteren Vektorraum ist $f$ genau dann linear, wenn die Verknüpfung $f \circ g : U \to W$ linear ist.}

Sind $g$ und $f$ linear, so wissen wir bereits, dass die Verknüpfung $f \circ g$ ebenfalls linear ist. Zu zeigen ist also nur, dass $f$ linear ist, wenn $f \circ g$ linear ist. 

\newpar

Sei also $f \circ g : U \to W$ linear. Seien $\vv_1, \vv_2 \in V$. 
Da $g$ surjektiv ist, existieren $\vu_1, \vu_2 \in U$ mit $g(\vu_1) = \vv_1$ und $g(\vu_2) = \vv_2$. Unter Anwendung der Linearität von $g$ und $f \circ g$ folgt:
\begin{align*}
	f(\vv_1 + \vv_2) 
	&= f(g(\vu_1) + (\vu_2))\\
	&= f(g(\vu_1 + \vu_2))\\
	&= (f \circ g)(\vu_1 + \vu_2)\\
	&= (f \circ g)(\vu_1) + (f \circ g)(\vu_2)\\
	&= f(g(\vu_1)) + f(g(\vu_2))\\
	&= f(\vv_1) + f(\vv_2)
\end{align*}
und:
\begin{align*}
	f(\lambda \vv_1)
	&=
	f(\lambda g(\vu_1))\\
	&=
	f(g(\lambda\vu_1))\\
	&=
	(f \circ g)(\lambda \vu_1)\\
	&=
	\lambda (f \circ g)(\vu_1)\\
	&=
	\lambda f(g(\vu_1))\\
	&=
	\lambda f(\vv_1)
\end{align*}
somit ist auch $f$ linear.

\clearpage

\section*{Übung 4.2}
\separator

\tit{Man zeige: Gegeben ein Vektorraum $V$ und eine lineare Abbildung $f : V \to V$ bildet die Menge der von $f$ festgehaltenen Vektoren alias aller \tib{Fixvektoren von $f$}
\begin{align*}
	V^f := \lr{v \in V \mid f(v) = v}
\end{align*}
stets einen Untervektorraum $V^f \subseteq V$.}

Zu zeigen ist also, dass $V^f$ den Nullvektor enthält und unter Vektoraddition und Skalarmultiplikation abgeschlossen ist.
\begin{enumerate}
	\item Es gilt $\vzero \in V^f$, da $f(\vzero) = \vzero$ sowieso für alle linearen Abbildungen $f$ gilt.
	\item Sind $\vv, \vw \in V^f$, so gilt mit der Linearität von $f$
	\begin{align*}
		f(\vv + \vw) 
		&=
		f(\vv) + f(\vw)
		=
		\vv + \vw,
	\end{align*}
	also $\vv + \vw \in V^f$
	\item Sind $\vv \in V^f$ und $\lambda \in K$, so gilt analog
	\begin{align*}
		f(\lambda \vv)
		&=
		\lambda f(\vv)
		= 
		\lambda \vv,
	\end{align*}
	also auch $\lambda \vv \in V$.
\end{enumerate}
Somit ist $V^f$ ein Untervektorraum von $V$.
\newpar
Diese Aufgabe ist eine Vorbereitung auf den späteren Umgang mit \tib{Eigenvektoren}, \tib{Eigenwerten} und \tib{Eigenräumen}. Ist $f : V \to V$ linear und $\lambda \in K$, nennen wir die Menge
\begin{align*}
	\Eig(f, \lambda) = \lr{v \in V \mid f(v) = \lambda v}
\end{align*}
den \tib{Eigenraum} von $f$ zum \tib{Eigenwert} $\lambda$, und dessen Elemente \tib{Eigenvektoren}. Mit den selben Argumenten wie in dieser Aufgabe folgt, dass auch $\Eig(f, \lambda)$ immer ein Untervektorraum von $V$ ist. 
\newpar
Eigenvektoren haben zahlreiche Anwendungen, sowohl in der reinen Mathematik als auch in praktischen Problemen der Numerik und der Natur- und Ingenieurwissenschaften, und werden deshalb noch zu einem sehr zentralen Teil der Vorlesungen LA I und II werden.

\clearpage

\section*{Übung 4.3}
\separator

\tit{Man zeige: Eine Abbildung $f : V \to W$ von Vektorräumen ist genau dann linear, wenn ihr Graph 
\begin{align*}
	\Gamma(f) = \lr{(v,f(v)) \mid v \in V} \subseteq V \times W
\end{align*}
ein Untervektorraum des Produkts ist.}

\begin{enumerate}
	\item "$\Longrightarrow$": Angenommen, $f$ ist linear. Seien $v_1, v_2 \in V$ und $\lambda \in K$. Dann gilt:
	\begin{enumerate}
		\item $f(\vzero_V) = \vzero_W$, also
		\begin{align*}
			\vzero_{V \times W} = (\vzero_V, \vzero_W) = (\vzero_V, f(\vzero_V)) \in \Gamma(f)
		\end{align*}
		\item $f(\vv_1 + \vv_2) = f(\vv_1) + f(\vv_2)$, also
		\begin{align*}
			(\vv_1, f(\vv_1)) + (\vv_2, f(\vv_2))
			=
			(\vv_1 + \vv_2, f(\vv_1) + f(\vv_2))
			=
			(\vv_1 + \vv_2, f(\vv_1 + \vv_2)) \in \Gamma(f)
		\end{align*}
		\item $f(\lambda \vv_1) = \lambda f(\vv_1)$, also
		\begin{align*}
			\lambda(\vv_1, f(\vv_1))
			=
			(\lambda \vv_1, \lambda f(\vv_1))
			=
			(\lambda \vv_1, f(\lambda \vv_1))
			\in
			\Gamma(f)
		\end{align*}
	\end{enumerate}
	Somit enthält $\Gamma(f)$ den Nullvektor und ist unter Addition und Skalarmultiplikation abgeschlossen, ist also ein Untervektorraum.
	\item "$\Longleftarrow$": Da $(\vv, f(\vv)) \in \Gamma(f)$ das einzige Element von $\Gamma(f)$ ist, deren erste Komponente $\vv$ ist folgt aus $(\vv, \vw) \in \Gamma(f)$ bereits $\vw = f(\vv)$. Ist $\Gamma(f)$ also ein Untervektorraum, so gilt:
	\begin{enumerate}
		\item Aus $(\vzero, \vzero) \in \Gamma(f)$ folgt $f(\vzero) = \vzero$.
		\item Aus $$(\vv_1, f(\vv_1)) + (\vv_2, f(\vv_2)) = (\vv_1 + \vv_2, f(\vv_1) + f(\vv_2)) \in \Gamma(f)$$ folgt $f(\vv_1 + \vv_2) = f(\vv_1) + f(\vv_2)$.
		\item Aus 
		\begin{align*}
			\lambda(\vv_1, f(\vv_1))
			=
			(\lambda \vv_1, \lambda f(\vv_1)) \in \Gamma(f)
		\end{align*}
		folgt $f(\lambda \vv_1) = \lambda f(\vv_1)$.
	\end{enumerate}
	Also ist $f$ linear.
\end{enumerate}
\end{document}
